\begin{proof}

\begin{enumerate}[label={$\langle1\rangle$\arabic*.},leftmargin=*]
\item Use the notation of P4R.7 and P4R.9.
The actual resolved pole support is
a section $H$ with rational trees,
and its inclusion in the ruled
resolution $T$ induces an isomorphism
on fundamental groups. Thus
$\pi_1(S)=\pi_1(C)$, where
$C=\operatorname{Sing}(S)_{\rm red}$.
The virtually abelian curve-group
calculation gives $e(C)\ge0$, and
P4.13 gives $e(S_{\rm reg})\ge0$.
Equality makes both numbers zero.
The same connected-curve calculation
shows that $C$ has exactly one
normalized component, of genus one,
and has no point-identification
cycles.

\item All rational components of the
resolved pole support are contracted
to points: a nonconstant proper map
from a rational curve to the
elliptic normalization of $C$ is
impossible. Hence the strict
conductor consists only of $H$.
The generic cusp has one
normalization branch, so
$H\to C$ has generic degree one
and is the normalization of $C$.
In particular $H$ and the ruling
base both have genus one.

\item Every exceptional curve of
$T\to N$ is now rational and
vertical for the ruling, since
$H$ is not contracted. The ruling
$r:T\to H$ descends to a holomorphic
map $\bar r:N\to H$. For an explicit
descent, the image of
$(\rho_N,r):T\to N\times H$
projects properly and bijectively
to $N$, because $r$ is constant
on each exceptional fibre. Its
normalized graph projects finitely
and birationally to the normal
space $N$, hence is isomorphic
to $N$.

\item The normalization singularities
are rational. Indeed a sequence
of point blowups of a
$\mathbb P^1$-bundle has
$R^1r_*\mathcal O_T=0$.
For a blowup this follows from
the two affine charts, or from
the exact sequence of the
exceptional $\mathbb P^1$ with
$H^1(\mathcal O_{\mathbb P^1})=0$;
the assertion for the whole
sequence follows by Leray.
Since $\rho_{N*}\mathcal O_T
=\mathcal O_N$, the low-degree
Leray sequence for
$r=\bar r\circ\rho_N$ gives
\[
 \bar r_*R^1\rho_{N*}\mathcal O_T=0.
\]
Here $R^2\bar r_*\mathcal O_N=0$
because $\bar r$ has one-dimensional
fibres. The sheaf
$R^1\rho_{N*}\mathcal O_T$ is
supported at finitely many
normalization singular points,
and direct image of a nonzero
such sheaf to $H$ is nonzero.
Thus it vanishes. In particular
$\chi(\mathcal O_N)
=\chi(\mathcal O_T)=0$.

\item We next check the scheme structure
of the conductor, to prove that
its image curve is smooth.
Write $I$ for the conductor ideal
on $S$ and on $N$, and let
$C_{\rm sch}\subset S$ and
$D_{\rm sch}\subset N$ be its
two zero schemes. Locally
$\mathcal O_S$ is Gorenstein
of dimension two and
$\nu_*\mathcal O_N$ is a
maximal Cohen--Macaulay module.
Finite duality identifies
its dual module with $I$,
which is again maximal
Cohen--Macaulay. The depth
lemma in
$0\to I\to\mathcal O_S
\to\mathcal O_{C_{\rm sch}}\to0$
therefore gives depth one
for the conductor curve.
At a generic cusp its local
length in the transverse
direction is one. A
one-dimensional Cohen--Macaulay
scheme which is generically
reduced is reduced: any
nilpotent ideal would have
zero-dimensional support,
contradicting the absence of
embedded associated points.
Consequently $C_{\rm sch}=C$.

\item Put $L=\nu^*\omega_S$, a
torsion line bundle by P4.12.
Finite duality on $N$ gives
\[
 I=\omega_N\otimes L^{-1}.
\]
Serre duality on the compact
Cohen--Macaulay surface $N$
gives $\chi(I)=\chi(L)$.
The rational-singularity
comparison just proved and
the projection formula give
$\chi(L)=\chi(\rho_N^*L)$.
Riemann--Roch on the smooth
surface $T$, with the rational
Chern class of this torsion
line bundle equal to zero,
gives $\chi(\rho_N^*L)
=\chi(\mathcal O_T)=0$.
Therefore
$\chi(\mathcal O_{D_{\rm sch}})
=\chi(\mathcal O_N)-\chi(I)=0$.

\item The exact conductor square is
\[
 0\longrightarrow\mathcal O_S
 \longrightarrow
 \nu_*\mathcal O_N\oplus\mathcal O_C
 \longrightarrow\nu_*\mathcal O_{D_{\rm sch}}
 \longrightarrow0.
\]
Taking Euler characteristics and
using $\chi(\mathcal O_S)=0$
from P4.12 yields
$\chi(\mathcal O_C)=0$.
Its normalization is the
smooth elliptic curve $H$,
so the normalization sequence
gives $\chi(\mathcal O_C)
 =-\sum_p\delta(C,p)$.
All delta invariants vanish;
thus $C$ is smooth.

\item Every point of $C$ now satisfies
P4.28, so $S$ has the product
cusp germ there and $N$ is
smooth along the conductor.
It was already smooth away
from the conductor. P4.22
therefore applies and makes
this elliptic ruled
normalization minimal.

\end{enumerate}
\end{proof}
